% GENERATED FILE. Edit canonical lessons, then run npm run generate:books.
% canonical-source-hash: fnv1a64:388f8168d669c633
% canonical-lessons: GE-C74-saft, GE-C74-wurst, GE-C74-schinken, GE-C74-salat, GE-C74-nudeln

\chapter{Juice, Sausage, Ham, Salad, Noodles}
\label{ch:ge-juice-sausage-ham-salad-noodles}
\hlchaptermodality{74}

\begin{chapteropening}
\textbf{By the end of this chapter:} \emph{I can say juice, a sausage, ham, a salad and noodles, pasta.}

Complete the last lesson of chapter 74: i can say juice, a sausage, ham, a salad and noodles, pasta.
\end{chapteropening}

\section[der Saft]{der Saft — juice}

\label{lesson:GE-C74-saft}



\begin{quote}
Before the new one: say the German for salt, then the German for sugar.
\end{quote}

\subsection*{You'll want to know: der Saft}
\textbf{der Saft} — \textquotedblleft{}juice\textquotedblright{}.

\begin{grammarlens}[title={What you've built}]
One more word for this chapter.
\end{grammarlens}

\subsection*{Your turn}
\begin{itemize}
\raggedright
  \item \emph{Say it:} \emph{der Saft}
  \item \emph{Say it:} \emph{der Saft}, once more
  \item \emph{Say it:} say \emph{der Saft}
  \item \emph{Recall:} say the German for salt, then the German for sugar, then say \emph{der Saft} again
\end{itemize}

\begin{tcolorbox}[breakable,colback=teal!4,colframe=teal!35!black,title={Before you move on}]
What does der Saft mean? (\textquotedblleft{}Juice\textquotedblright{}.) Say it once more.
\end{tcolorbox}
\section[die Wurst]{die Wurst — a sausage}

\label{lesson:GE-C74-wurst}



\begin{quote}
Before the new one: say the German for sugar, then the German for juice.
\end{quote}

\subsection*{You'll want to know: die Wurst}
\textbf{die Wurst} — \textquotedblleft{}a sausage\textquotedblright{}.

\begin{grammarlens}[title={What you've built}]
One more word for this chapter.
\end{grammarlens}

\subsection*{Your turn}
\begin{itemize}
\raggedright
  \item \emph{Say it:} \emph{die Wurst}
  \item \emph{Say it:} \emph{die Wurst}, once more
  \item \emph{Say it:} say \emph{die Wurst}
  \item \emph{Recall:} say the German for sugar, then the German for juice, then say \emph{die Wurst} again
\end{itemize}

\begin{tcolorbox}[breakable,colback=teal!4,colframe=teal!35!black,title={Before you move on}]
What does die Wurst mean? (\textquotedblleft{}A sausage\textquotedblright{}.) Say it once more.
\end{tcolorbox}
\section[der Schinken]{der Schinken — ham}

\label{lesson:GE-C74-schinken}



\begin{quote}
Before the new one: say the German for juice, then the German for a sausage.
\end{quote}

\subsection*{You'll want to know: der Schinken}
\textbf{der Schinken} — \textquotedblleft{}ham\textquotedblright{}.

\begin{grammarlens}[title={What you've built}]
One more word for this chapter.
\end{grammarlens}

\subsection*{Your turn}
\begin{itemize}
\raggedright
  \item \emph{Say it:} \emph{der Schinken}
  \item \emph{Say it:} \emph{der Schinken}, once more
  \item \emph{Say it:} say \emph{der Schinken}
  \item \emph{Recall:} say the German for juice, then the German for a sausage, then say \emph{der Schinken} again
\end{itemize}

\begin{tcolorbox}[breakable,colback=teal!4,colframe=teal!35!black,title={Before you move on}]
What does der Schinken mean? (\textquotedblleft{}Ham\textquotedblright{}.) Say it once more.
\end{tcolorbox}
\section[der Salat]{der Salat — a salad; lettuce}

\label{lesson:GE-C74-salat}



\begin{quote}
Before the new one: say the German for a sausage, then the German for ham.
\end{quote}

\subsection*{You'll want to know: der Salat}
\textbf{der Salat} — \textquotedblleft{}a salad; lettuce\textquotedblright{}.

\begin{grammarlens}[title={What you've built}]
One more word for this chapter.
\end{grammarlens}

\subsection*{Your turn}
\begin{itemize}
\raggedright
  \item \emph{Say it:} \emph{der Salat}
  \item \emph{Say it:} \emph{der Salat}, once more
  \item \emph{Say it:} say \emph{der Salat}
  \item \emph{Recall:} say the German for a sausage, then the German for ham, then say \emph{der Salat} again
\end{itemize}

\begin{tcolorbox}[breakable,colback=teal!4,colframe=teal!35!black,title={Before you move on}]
What does der Salat mean? (\textquotedblleft{}A salad; lettuce\textquotedblright{}.) Say it once more.
\end{tcolorbox}
\section[die Nudeln]{die Nudeln — noodles, pasta}

\label{lesson:GE-C74-nudeln}



\begin{quote}
Before the new one: say the German for ham, then the German for a salad; lettuce.
\end{quote}

\subsection*{You'll want to know: die Nudeln}
\textbf{die Nudeln} — \textquotedblleft{}noodles, pasta\textquotedblright{}.

\begin{grammarlens}[title={What you've built}]
One more word for this chapter.
\end{grammarlens}

\subsection*{Your turn}
\begin{itemize}
\raggedright
  \item \emph{Say it:} \emph{die Nudeln}
  \item \emph{Say it:} \emph{die Nudeln}, once more
  \item \emph{Say it:} say \emph{die Nudeln}
  \item \emph{Recall:} say the German for ham, then the German for a salad; lettuce, then say \emph{die Nudeln} again
\end{itemize}

\begin{tcolorbox}[breakable,colback=teal!4,colframe=teal!35!black,title={Before you move on}]
What does die Nudeln mean? (\textquotedblleft{}Noodles, pasta\textquotedblright{}.) Say it once more.
\end{tcolorbox}
